% sections/methods.tex -- fixed settings typed after checking docs/protocol.md, the configs and the code; measured values
% are macros. Procedural detail is in sections/supp_methods.tex (S1 Appendix, part A).
\subsection*{Study design}
All models were trained and evaluated with one protocol (Fig~\ref{fig:pipeline}). Four CNNs formed a 2 $\times$ 2 design in which attention gates and multi-scale blocks were switched on or off on a fixed
U-Net backbone. SAM was fine-tuned in two configurations: on whole images reduced to 1024 $\times$ 1024 and on four
full-resolution tiles, with an inference-only variant that used overlapping tiles. Each trained model had three seeds,
and all predictions were scored at 2048 $\times$ 2048 inside a field-of-view (FOV) mask at a threshold chosen
on the validation set. Zero-shot SAM with point prompts from the ground truth was a separate experiment.

\begin{figure}[!h]
\figpreview{figures/fig_pipeline.pdf}
\caption{\CapPipeline}
\label{fig:pipeline}
\end{figure}

\subsection*{Data, audit and split}
This is a retrospective analysis of a public dataset. FIVES~\cite{jin2022fives} contains colour fundus photographs of
2048 $\times$ 2048 pixels taken with a Topcon TRC-NW8 camera (50$^\circ$ field of view) between 2016 and 2021 at one
hospital, the Second Affiliated Hospital of Zhejiang University, from 573 patients (mean age 48 years, standard
deviation 19.87; 469 images from female and 331 from male subjects; all Asian), with 200 images each of age-related
macular degeneration (AMD), diabetic retinopathy (DR), glaucoma and normal eyes. Every image has a pixelwise vessel
annotation made by trained medical staff and verified by ophthalmologists, which we used as the reference standard without
editing~\cite{jin2022fives}. We used the copy on the Hugging Face hub (\texttt{tyluan/FIVES}; CC BY 4.0~\cite{jin2022fivesdata});
all images were RGB and all masks binary. The
official split has \NTrainOfficial{} training and \NTestOfficial{} test images (\NTestPerCategory{} per category); the
sample size was fixed by the dataset. No other dataset was used.

Before any model was trained, we checked all images and masks of this copy for exact and near duplicates, empty masks,
annotations outside the eye and identifying marks (Table~\ref{tab:data}, Fig~\ref{fig:audit}; procedures and parameters in
\nameref{si:appendix}). A later check showed that the copy is identical to the official release on figshare, so the
findings apply to FIVES itself. Many images carry burned-in text outside the eye that may contain fragments of patient
names; it was compared only as binary patterns and never read (Ethics statement). Three training images were excluded: two very dark images with empty masks, and one
showing the same eye as a test image.

FIVES provides no patient identifiers. To keep duplicate and possibly same-patient images together, we grouped
byte-identical files (\NGroupPairsIdentical{} pairs), images with an identical marking pattern and consecutive numbers
(\NGroupPairsAdjacent{} pairs), and isolated pairs of images with an identical marking pattern (\NGroupPairsIsolated{}
pairs): \NGroups{} groups containing \NGroupedImages{} images. The remaining \NTrainVal{} official training images were
split into \NTrain{} training and \NVal{} validation images, \NValPerCategory{} per category, with groups assigned whole
(seed \SplitSeed{}); no group crosses the split. The official test set was not changed.

% tables/data.tex -- generated by paper/scripts/make_assets.py. Do not edit by hand.
\begin{table}[!ht]
\begin{adjustwidth}{-0.5in}{-0.5in}
\centering
\caption{{\bf Data, split and audit.} Top: images per disease category in the official FIVES split and in the split used here (training and validation drawn from the official training images with seed \SplitSeed{}, \NValPerCategory{} validation images per category; \NGroups{} groups of duplicate or possibly same-patient images, \NGroupedImages{} images, never split). Bottom: findings of the data audit of all \NImages{} images and masks in the copy used (Hugging Face mirror of FIVES, identical to the official release), made before any model was trained. Markings were detected and compared as binary patterns only and were never read.}
\begin{tabular}{p{1.45in}cccccc}
\hline
{\bf Category} & {\bf Official train} & {\bf Official test} & {\bf Train} & {\bf Validation} & {\bf Excluded} & {\bf Test} \\ \thickhline
AMD (A) & 150 & 50 & 120 & 30 & 0 & 50 \\ \hline
DR (D) & 150 & 50 & 119 & 30 & 1 & 50 \\ \hline
Glaucoma (G) & 150 & 50 & 118 & 30 & 2 & 50 \\ \hline
Normal (N) & 150 & 50 & 120 & 30 & 0 & 50 \\ \hline
Total & 600 & 200 & 477 & 120 & 3 & 200 \\ \hline
\multicolumn{7}{l}{{\bf Audit finding}} \\ \thickhline
Same image file, two different masks & \multicolumn{3}{p{2.0in}}{\NDupPairsTrain{} pairs in train, \NDupPairsTest{} in test; mask Dice over the whole image \DupTrainMaskDiceMin{}--\DupTrainMaskDiceMax{} (train), \DupTestMaskDice{} (test)} & \multicolumn{3}{p{1.9in}}{Train pairs kept on one side of the split; test pair counted as two images} \\ \hline
Empty ground-truth mask & \multicolumn{3}{p{2.0in}}{\NEmptyMasks{} very dark images (a near-identical pair)} & \multicolumn{3}{p{1.9in}}{Excluded} \\ \hline
Same eye in official train and test & \multicolumn{3}{p{2.0in}}{\NSameEyePairs{} pair, vessel-layout IoU \SameEyeIoU{}} & \multicolumn{3}{p{1.9in}}{Training image excluded; test set unchanged} \\ \hline
Annotation lines outside the eye & \multicolumn{3}{p{2.0in}}{\NBorderMasks{} masks, \NBorderPx{} pixels} & \multicolumn{3}{p{1.9in}}{All scoring inside a fitted FOV mask; masks not edited} \\ \hline
Burned-in text markings outside the eye & \multicolumn{3}{p{2.0in}}{\NMarkings{} of \NImages{} images; possibly identifying} & \multicolumn{3}{p{1.9in}}{Never read; pixelated in all figures; compared only as binary patterns} \\ \hline
Possibly same patient (identical marking pattern, isolated pair) & \multicolumn{3}{p{2.0in}}{\NIsoPairs{} pairs: \NIsoTrainTest{} train--test, \NIsoValTest{} validation--test, \NIsoTrainVal{} train--validation, \NIsoTestTest{} within test} & \multicolumn{3}{p{1.9in}}{Grouped within the official training set; pairs involving test remain} \\ \hline
\end{tabular}
\label{tab:data}
\end{adjustwidth}
\end{table}


\begin{figure}[!h]
\figpreview{figures/fig_audit_raster.png}
\caption{\CapAudit}
\label{fig:audit}
\end{figure}

\subsection*{Preprocessing and training targets}
Contrast was enhanced with CLAHE (contrast-limited adaptive histogram equalization)~\cite{zuiderveld1994clahe} on the
lightness channel of the LAB colour space at full resolution (clip limit 2.0, 8 $\times$ 8 tiles), keeping three channels.
For the CNNs and whole-image SAM the enhanced images were reduced by area averaging to 512 $\times$ 512 and
1024 $\times$ 1024. CNN inputs were normalized with the per-channel mean and standard deviation of the training images;
SAM inputs used only SAM's own normalization. Training images and masks were flipped horizontally and vertically (each with
probability 0.5) and rotated by a random multiple of 90$^\circ$.

Reducing a mask to the training resolution either removes thin vessels or widens them. In a comparison on 30 training
masks made before any model was trained (\nameref{si:resize}), block averaging without a threshold kept the most
information (Dice \ResizeFiveSoftDice{} at 512, against \ResizeFiveAreaHalfDice{} with a threshold of 0.5) without widening
vessels (precision \ResizeFiveSoftPrec, against \ResizeFiveMaxpoolPrec{} for max-pooling). We therefore trained the CNNs on
these soft masks at 512 and whole-image SAM on soft masks at 1024. In a soft mask a thin vessel is stored as a value below
0.5, so its recovery depends on the threshold (thin-vessel recall of the soft 512 mask: \ResizeFiveSoftThin{} at 0.5,
\ResizeFiveSoftQuarterThin{} at 0.25). Tiled SAM was trained on full-resolution tiles, whose masks are binary.

The CNNs work at 512 $\times$ 512. A pilot with the U-Net at 1024 $\times$ 1024 used a rule fixed before its result was seen:
all CNNs would move to 1024 only if validation thin-vessel recall rose by at least 0.05 while validation Dice fell by at most
0.005. Thin-vessel recall rose by \PilotThinGain{} (Dice \PilotDiceGain), so the CNNs stayed at 512 (\nameref{si:pilot}).

\subsection*{Models}
\emph{CNNs.} The baseline U-Net~\cite{ronneberger2015unet} has four down-sampling and four up-sampling levels with 64 to
1024 channels; each block applies two 3 $\times$ 3 convolutions with batch normalization and ReLU, followed by
2 $\times$ 2 max-pooling in the encoder and 2 $\times$ 2 transposed convolution in the decoder, and a 1 $\times$ 1
convolution produces the logits (\UNetParams{} parameters). The variants keep this backbone (Table~\ref{tab:models}).
Additive \emph{attention gates}~\cite{oktay2018attention} multiply each of the four skip connections by coefficients
computed from the skip features and the decoder features one level coarser (definition in \nameref{si:appendix}, A10).
The \emph{multi-scale block} replaces the second 3 $\times$ 3 convolution
of all nine blocks with three parallel 3 $\times$ 3 convolutions with dilation 1, 2 and 4~\cite{yu2016dilated}, whose
output channels are split rather than tripled (for 64 channels: 22, 21 and 21), concatenated and fused by batch
normalization, ReLU, a 1 $\times$ 1 convolution, batch normalization and ReLU. Attention gates add \AttParamsPct\% to the
parameters, multi-scale blocks \MSParamsPct\% and both \MSAttParamsPct\%. With the same seed, the Attention U-Net starts
from the U-Net's backbone weights.

% tables/models.tex -- generated by paper/scripts/make_assets.py. Do not edit by hand.
\begin{table}[!ht]
\begin{adjustwidth}{-0.5in}{-0.5in}
\centering
\caption{{\bf Models compared.} Everything not listed is shared by all trained models: data split, CLAHE preprocessing, augmentation (flips and 90\textdegree{} rotations), loss (binary cross-entropy + soft Dice), Adam with learning rate 10\textsuperscript{-4}, effective batch 4, early stopping on validation Dice at 2048 (patience 20, at most 150 epochs), mixed precision, three seeds, threshold chosen on validation, and scoring at 2048 inside the FOV. Parameters in millions (trainable / total).}
\begin{tabular}{p{1.2in}p{1.25in}p{1.2in}p{0.9in}p{0.9in}p{0.7in}}
\hline
{\bf Model} & {\bf Change vs U-Net} & {\bf Network input} & {\bf Training target} & {\bf Output grid} & {\bf Params (M)} \\ \thickhline
U-Net & -- & CLAHE image, area-resized to 512 & soft mask at 512 & 512 & \UNetParamsTrainableM{} / \UNetParamsM{} \\ \hline
Attention U-Net & attention gates on 4 skips & as U-Net & as U-Net & 512 & \AttParamsTrainableM{} / \AttParamsM{} \\ \hline
Multi-Scale U-Net & multi-scale block in 9 blocks & as U-Net & as U-Net & 512 & \MSParamsTrainableM{} / \MSParamsM{} \\ \hline
Multi-Scale Attention U-Net & both & as U-Net & as U-Net & 512 & \MSAttParamsTrainableM{} / \MSAttParamsM{} \\ \hline
SAM, tiled & SAM ViT-B, LoRA r=4 + decoder & 4 native 1024 crops & binary native crop & 512 (4$\times$256 stitched) & \SAMTiledParamsTrainableM{} / \SAMTiledParamsM{} \\ \hline
SAM, overlapping tiles & same models as tiled & 9 native 1024 crops, stride 512 & -- & 512 (tent-blended) & \SAMOverlapParamsTrainableM{} / \SAMOverlapParamsM{} \\ \hline
SAM, whole image & SAM ViT-B, LoRA r=4 + decoder & CLAHE image, area-resized to 1024 & soft mask at 1024 & 1024 (from 256) & \SAMWholeParamsTrainableM{} / \SAMWholeParamsM{} \\ \hline
\end{tabular}
\label{tab:models}
\end{adjustwidth}
\end{table}


\emph{Fine-tuned SAM.} We started from the released weights of SAM with the ViT-B image
encoder~\cite{kirillov2023sam,dosovitskiy2021vit} (\OracleParams{} parameters). Following the LoRA placement of SAMed~\cite{zhang2023samed}, rank-4 LoRA~\cite{hu2022lora} was added to the query and value projections of all 12 encoder blocks; the LoRA layers and the
mask decoder were trained (\SAMTiledParamsTrainable{} parameters) and everything else was frozen. Unlike SAMed, we kept the
prompt encoder frozen and gave no prompts (implementation and gradient
check in \nameref{si:appendix}). \emph{Whole-image SAM} receives the enhanced image at 1024 $\times$ 1024; its 256 $\times$ 256
decoder logits are enlarged to 1024 by SAM's post-processing and compared with the soft 1024 masks. \emph{Tiled SAM}
receives the four non-overlapping 1024 $\times$ 1024 tiles of the full-resolution image; each tile's logits, enlarged to
1024, are compared with its binary mask, and at inference the four 256 $\times$ 256 logit grids are joined into one
512 $\times$ 512 grid, the grid of the CNNs. A tiled epoch covers the same image content as a CNN epoch in four times as
many optimizer steps.

On the validation set, the non-overlapping tiles left a discontinuity along the tile borders (\nameref{si:seam}). We kept
the tiled models unchanged as the main SAM and added inference with \emph{overlapping tiles}: nine 1024 $\times$ 1024
tiles at a stride of 512 pixels, whose logits are combined as a weighted average with triangular weights, normalized by
the sum of the weights at each pixel.

\emph{Zero-shot SAM with oracle prompts.} As a separate experiment, the unmodified SAM received $K$ positive and $K$
negative points sampled from each image's own ground truth ($K \in \{5, 10, 20\}$ and threshold chosen on validation;
\nameref{si:appendix}). Because the prompts use every image's ground truth, it is reported separately
(\nameref{si:oracle}).

\subsection*{Training}
The loss was the sum of binary cross-entropy on the logits and a soft Dice loss $1 - (2\sum p t + 1)/(\sum p + \sum t + 1)$
per image on the sigmoid outputs $p$ and targets $t$, computed in single precision and applied to soft targets as they
are. Optimization used Adam~\cite{kingma2015adam} with a learning rate of $10^{-4}$ and batches of four images (SAM: two
accumulated micro-batches of two) with mixed precision~\cite{micikevicius2018mixed}, for at most 150 epochs, stopping after
20 epochs without improvement of the validation Dice and keeping the best checkpoint. Validation Dice was computed as the
test data are scored: maps enlarged to 2048 $\times$ 2048, threshold 0.5, original masks, inside the FOV. Each model was
trained with seeds 0, 1 and 2, which set the initial weights, the image order and the augmentation, on one NVIDIA L40S
GPU per run (\nameref{si:runs}); seeds are reported separately, and no ensembling was used. A check before training
showed that this loss drives soft targets towards 0 or 1: at its optimum a target of 0.25, the value of a one-pixel vessel at 512, becomes about \LossOptQuarter{} (\nameref{si:appendix}).

\subsection*{Evaluation}
Every metric was computed per image inside an FOV mask: the union of a circle fitted to the eye boundary and the bright eye
region (\nameref{si:appendix}), which includes dim rims and excludes the burned-in text. The FOV covered
\FOVShareMin\%--\FOVShareMax\% of an image. Probability maps were enlarged bilinearly to
2048 $\times$ 2048 and thresholded. For each run the threshold was chosen on the validation set from 0.05 to 0.95 in steps of
0.05 as the value with the highest mean Dice (ties: the lowest) and applied to that run's test predictions. As a secondary
analysis decided before test evaluation, every run was also scored at a fixed threshold of 0.5, and thin-vessel recall was
recorded at every grid threshold.

With $P$ and $G$ the predicted and ground-truth vessel sets inside the FOV, we report Dice, intersection over union (IoU),
precision and recall,
\begin{equation}
\mathrm{Dice} = \frac{2\lvert P \cap G\rvert}{\lvert P\rvert + \lvert G\rvert},\quad
\mathrm{IoU} = \frac{\lvert P \cap G\rvert}{\lvert P \cup G\rvert},\quad
\mathrm{precision} = \frac{\lvert P \cap G\rvert}{\lvert P\rvert},\quad
\mathrm{recall} = \frac{\lvert P \cap G\rvert}{\lvert G\rvert}.
\end{equation}
For thin vessels we used centreline recall by vessel width. The ground truth was skeletonized~\cite{vanderwalt2014skimage,zhang1984thinning}
and each skeleton point $s$ received the local radius $r(s)$, the Euclidean distance transform of the ground truth at
$s$~\cite{virtanen2020scipy}. Points with $r < 3$ pixels were thin, $3 \le r \le 7$ medium and $r > 7$ thick (the cut-offs
of a recent preprint~\cite{amarnath2026resolution}). In the training masks \ShareThin\% of skeleton points were thin, \ShareMed\% medium and \ShareThick\% thick. The recall
of class $c$ is the fraction of its skeleton points $S_c$ inside the prediction dilated by one pixel ($3 \times 3$ square
$B_1$),
\begin{equation}
R_c = \frac{\lvert \{ s \in S_c : s \in P \oplus B_1 \} \rvert}{\lvert S_c \rvert};
\end{equation}
as a sensitivity check the classes were also formed with cut-offs of 4 and 8 pixels. With $S_P$ and $S_G$ the skeletons of
the two masks, clDice~\cite{shit2021cldice} is
\begin{equation}
\mathrm{clDice} = \frac{2\,T_{\mathrm{prec}}\,T_{\mathrm{sens}}}{T_{\mathrm{prec}} + T_{\mathrm{sens}}},\quad
T_{\mathrm{prec}} = \frac{\lvert S_P \cap G\rvert}{\lvert S_P\rvert},\quad
T_{\mathrm{sens}} = \frac{\lvert S_G \cap P\rvert}{\lvert S_G\rvert}.
\end{equation}
Boundary accuracy was measured with normalized surface Dice (NSD)~\cite{nikolov2021nsd} in a boundary-pixel approximation,
with $\partial M$ a mask minus its $3 \times 3$ erosion and $(\partial M)_\tau$ the pixels within $\tau = 2$ pixels of it:
\begin{equation}
\mathrm{NSD} = \frac{\lvert \partial P \cap (\partial G)_\tau\rvert + \lvert \partial G \cap (\partial P)_\tau\rvert}
{\lvert \partial P\rvert + \lvert \partial G\rvert}.
\end{equation}
Skeletons, boundaries and radii were computed on the full masks and then restricted to the FOV, so the FOV border creates
no artificial vessel ends; a ratio with an empty denominator was 1 if both sets were empty and 0 otherwise. As
image-derived vessel measures we computed vessel density (vessel pixels per FOV pixel) and the fractal
dimension~\cite{liew2008fractal}, here by box counting with boxes of 2--512 pixels, for prediction and ground truth with
the same FOV.

\emph{Mask-resampling references.} Each ground truth was reduced to soft masks at 512 $\times$ 512 (the CNN and tiled-SAM
grid) and at 256 $\times$ 256 (the whole-image SAM decoder grid, enlarged to 1024 as SAM does), enlarged back to 2048 and
scored with the same threshold rule. These references measure what resizing alone loses. They are not upper bounds: a
model's logits are enlarged before the sigmoid and are not constrained to block averages.

\subsection*{Statistical analysis}
Model-level results are the mean $\pm$ standard deviation across the three seeds of each seed's mean over the
\NTestImages{} test images, and describe training-seed variability. We call a difference consistent across seeds when every
seed of one model lies on the same side of every seed of the other; with three seeds this describes the runs evaluated, not
general robustness. Complete separation of three runs from three is the most extreme of 20 equally likely arrangements, so it corresponds to
an exact one-sided permutation p-value of 0.05, and failing to reach it is weak evidence that an effect is absent. The seeds
of the CNNs are coupled (the same seed gives the same image order and augmentation, and the same backbone initialization
for a model and its attention variant); only the post hoc analyses treat them as paired. Image-level analyses averaged each image's
metric over a model's seeds first. Each automatic model was compared with the U-Net by a two-sided paired Wilcoxon signed-rank test~\cite{wilcoxon1945,virtanen2020scipy} with Holm
correction~\cite{holm1979} within each metric across the six comparisons. For the 2 $\times$ 2 design we computed per image
the attention main effect (mean of Attention U-Net $-$ U-Net and Multi-Scale Attention U-Net $-$ Multi-Scale U-Net), the
multi-scale main effect (analogous) and the interaction (Multi-Scale Attention U-Net $-$ Multi-Scale U-Net $-$ Attention
U-Net $+$ U-Net), with 95\% percentile bootstrap confidence intervals (CIs)~\cite{efron1979bootstrap} from 10,000 resamples of
the test images. Image-level tests and intervals do not include training-seed variability, so we report both kinds of
variability and give effect sizes rather than p-values alone. The threshold rules (validation-selected as primary, fixed 0.5 as secondary) were fixed before test evaluation; no single
primary metric was designated, thin-vessel recall being the focus named in the study objectives, and Holm correction was
applied within each metric only. Results by disease category are descriptive; vessel measures
were compared with the ground truth by the mean difference, Bland--Altman limits of agreement~\cite{bland1986} and Pearson
correlation. Reporting follows the CLAIM 2024 checklist~\cite{tejani2024claim,mongan2020claim} (\nameref{si:claim}).

\subsection*{Analysis decisions and later additions}
The split, preprocessing, targets, loss, stopping rule, FOV scoring, metric definitions, width cut-offs, threshold rule,
fixed-threshold secondary analysis, resolution rule and statistical analysis were decided and recorded before test
evaluation, in a public commit made before any test prediction (\texttt{0ee73608db87db9e0b6f2ff46d29f2faa75fd1dc}). Three
additions came later. The tile-seam check and the overlapping-tile variant were decided on validation data after that
commit but before any SAM test prediction. After the first test evaluation, the two SAM configurations, which until then had
one seed, were trained with seeds 1 and 2 using unchanged configurations, so that seed variability could be estimated as for
the CNNs, and overlapping-tile inference was run for them; the CNNs had three seeds from the start. As first observed, the
seed-0 tiled SAM exceeded the three-seed U-Net mean in thin-vessel recall by \SAMTiledSeedAThinMinusUNet. The
annotation-agreement analysis below was added during writing. Each configuration predicted the test set once, without
test-time augmentation or post-processing.

During revision, after the main test evaluation, we added post hoc analyses that re-scored saved probability maps or
re-analysed saved results, without new predictions or test-set choices (\nameref{si:appendix}, A9): width-stratified
centreline precision, classing each predicted centreline pixel by its own predicted width; thin-vessel recall at matched
precision, with each run's threshold set on validation to match the U-Net's validation precision; clDice and NSD at 0.5;
a thin class without skeleton spurs; bootstrap intervals resampling clusters of possibly same-patient test images and,
jointly, training seeds; and a comparison with the official release.

\subsection*{Agreement between two available annotations of the same image}
For the \AgreeNPairs{} pairs of byte-identical image files with different masks (\AgreeNTrain{} training, \AgreeNTest{}
test; all AMD), each mask was scored against the other with the metrics above in both directions and the two values were
averaged. This post hoc descriptive analysis uses labels only; as the dataset does not say who drew the masks, it measures
agreement between two available annotations, not inter-rater agreement.

\subsection*{Computational cost and software}
Inference time, operations and peak memory were measured for the seed-0 models on one L40S with one image per call over
the test images (details in \nameref{si:appendix}); training time and memory are those of the seed-0 runs. The analysis used
Python 3.11, PyTorch 2.14.1 (CUDA 12.6)~\cite{paszke2019pytorch}, segment-anything 1.0, NumPy 2.4.6, SciPy
1.17.1~\cite{virtanen2020scipy}, scikit-image 0.26.0~\cite{vanderwalt2014skimage}, OpenCV 5.1.0 and pandas 3.1.1.
